\chapter{Discussion}

\section{Interpretation of Findings}
The results support the thesis claim that educator-centered workflow design can deliver major usability gains without giving up standards-compatible output. Improvements are consistent with cognitive-load reduction principles and usability heuristics from prior work \citep{sweller1988,nielsen1994,clarkmayer2016}.

\section{Pedagogical Implications}
Interactive video authoring becomes feasible for instructors who are not CMS specialists. This rebalances effort from technical setup toward instructional design quality, improving the likelihood of sustained adoption.

\section{Architectural Implications}
A service-oriented backend with explicit validation boundaries proved effective for combining traditional CRUD operations with probabilistic AI outputs. The architecture remains extensible for advanced analytics, rubric alignment, and multilingual support.

\section{Limitations}
The current version has the following limitations:
\begin{itemize}[leftmargin=2em]
\item Evaluation scale is limited and may not represent all teaching contexts.
\item Mobile-first authoring remains partial.
\item Full coverage of all H5P content types is not yet implemented.
\item Long-running AI tasks require more robust queueing for very large deployments.
\end{itemize}

\section{Generalizability}
While rooted in the VNU-HCM context, the design principles and architecture patterns are broadly transferable to higher-education institutions with similar LMS and cloud constraints.
